% !TeX root = demo.tex

\chapter{内容示例}
\section{用户环境示例}

\begin{tcolorbox}
	\begin{description}
		\item [GRE备考指南] \url{https://qyxf.site/latest/GRE备考指南-v2.0.pdf}
		\item [军事理论教程] \url{https://qyxf.site/latest/军事理论教程.pdf}
		\item [分析力学笔记] \url{https://qyxf.site/latest/分析力学笔记-v1.0.pdf}
		\item [大学物理题解] \url{https://qyxf.site/latest/大物题解（上）.pdf}
		\item [实变函数习题解答] \url{https://qyxf.site/latest/实变函数习题解答.pdf}
		\item [计算方法撷英] \url{https://qyxf.site/latest/计算方法撷英-v1.1.pdf}
		\item [计算机程序设计指南] \url{https://qyxf.site/latest/计算机设计程序指南.pdf}
	\end{description}
\end{tcolorbox}

\begin{pequation}
	\textit{Euler-Lagrangian equation}
	$$
	\dd{}{t}\p{\lag}{\dvr_i}=
	$$
\end{pequation}

\begin{define}
	极限就是超越自我。
\end{define}

\begin{theorem}
	任何极限都可以直接观察得出。
\end{theorem}

\begin{lemma}
	以上内容，纯属扯淡。
\end{lemma}

\begin{corollary}
	这是一个推论。
\end{corollary}

\begin{note}
	好好学习，天天向上。
\end{note}

\begin{alert}
	今天你学习了吗？
\end{alert}

\section{列表样式}
\begin{itemize}
	\item 这是第一层
	\item 这也是第一层
	\begin{itemize}
		\item 这是第二层
		\begin{itemize}
			\item 这是第三层
		\end{itemize}
	\end{itemize}
\end{itemize}

\begin{enumerate}
	\item 这是第一层
	\item 这也是第一层
	\begin{enumerate}
		\item 这是第二层
		\begin{enumerate}
			\item 这是第三层
		\end{enumerate}
	\end{enumerate}
\end{enumerate}

\section{正文示例}

\textbf{微分学}（\emph{differential calculus}）是微积分的一部分，是通过\emph{导数}和\emph{微分}来研究曲线斜率、加速度、最大值和最小值的一门学科，也是探讨特定数量变化速率的学科。微分学是微积分的两个主要分支之一，另一个分支则是\textbf{积分学}，探讨曲线下的面积。


\begin{table}[htbp]
	\centering
	\caption{常用导数}
	\begin{tabular}{cccc}
		\toprule
		\textbf{原函数} & \textbf{导函数}   & \textbf{原函数} & \textbf{导函数} \\
		\midrule
		$C$             & $0$               & $\ln x$         & $\frac{1}{x}$   \\
		$x^\mu$         & $\mu x^{\mu - 1}$ & $\sin x$        & $\cos x$        \\
		$e^x$           & $e^x$             & $\cos x$        & $-\sin x$       \\
		\bottomrule
	\end{tabular}
\end{table}

……几乎所有量化的学科中都有\textbf{微分}的应用。例如在物理学中，运动物体其\emph{位移}对时间的导数即为其\emph{速度}，\emph{速度}对时间的导数就是\emph{加速度}、物体\emph{动量}对时间的导数即为物体所受的\emph{力}，重新整理后可以得到牛顿第二运动定律 $F=ma$ 。化学反应的\emph{化学反应速率}也是导数。在运筹学中，会透过导数决定在运输或是设计上最有效率的做法。

\begin{figure}[htbp]
	\centering
	\includegraphics[width=.3\textwidth]{qyxf-cover.pdf}
	\caption{V2 版本的封面图片}
	\label{fig:qyxf-logo}
\end{figure}

导数常用来找函数的\emph{极值}。含有微分项的方程式称为\textbf{微分方程}，是自然现象描述的基础。微分以及其广义概念出现在许多数学领域中，例如\emph{复分析}、\emph{泛函分析}、\emph{微分几何}、\emph{测度}及\emph{抽象代数}\footnote{以上内容摘自维基百科中文词条 --- 微分学：\url{https://zh.wikipedia.org/wiki/微分学}。}。

\section{引导命令示例}

\exercise{1} 试用配方法求解方程：
\begin{equation}\label{eq:quadratic}
	ax^2 + bx + c = 0
\end{equation}

\solve 首先，方程左右两侧同除以 $a$，得到
\[ x^2 + \frac bax + \frac ca = 0 \]
根据一次项来配方，按公式 $(x+A)^2=x^2+2Ax+A^2$ 配出常数项：
\[ x^2 + \frac bax + \left(\frac b{2a}\right)^2 + \frac ca - \left(\frac b{2a}\right)^2 = 0 \]
配方并移项得到
\[ \left(x + \frac b{2a}\right)^2 = \frac {b^2}{4a^2} - \frac ca \]
方程左右开方，得
\[ x + \frac b{2a} = \pm \sqrt{\frac {b^2}{4a^2} - \frac ca} \]
从而得到方程 \eqref{eq:quadratic} 之解为
\begin{equation}
	x = - \frac b{2a} \pm \sqrt{\frac {b^2}{4a^2} - \frac ca}
\end{equation}
该式即为一元二次方程的\textbf{通用求根公式}。


\analysis 在这一问题中，需要注意以下几点 \cite{texbook,latex}：
\begin{itemize}
	\item ……
	\item ……
	\item ……
\end{itemize}

%\begin{thebibliography}{99}
%  \bibitem{texbook} KNUTH~D~E. The \TeX book [M]. Addison-Wesley: Reading, 1986.
%  \bibitem{latex} 刘海洋. \LaTeX 入门 [M]. 人民邮电出版社: 北京, 2013.
%\end{thebibliography}